\paragraph{Poisson--Cauchy 位置混合的四阶熵主项、黄金离散化与桥接判据}
以下结论在同尺度位置混合子类上给出可审计的主常数与桥接条件，并与本节既有 Poisson 粗粒化框架保持一致。

\begin{theorem}[Poisson--Cauchy 位置混合的 KL 四阶主项与六阶余项]\label{thm:xi-poisson-cauchy-mixture-kl-t4-minimal-third}
\leanverified{paper\_xi\_poisson\_cauchy\_mixture\_kl\_t4\_minimal\_third}
设 \(\mu\in\mathcal P(\RR)\) 满足
\[
\int_{\RR} y\,\dd\mu(y)=0,\qquad
\sigma^2:=\int_{\RR} y^2\,\dd\mu(y),\qquad
M_6:=\int_{\RR}|y|^6\,\dd\mu(y)<\infty.
\]
对 \(t>0\) 定义
\[
P_t(x):=\frac{1}{\pi}\frac{t}{x^2+t^2},\qquad
p_t(x):=(P_t*\mu)(x),\qquad
q_t(x):=P_t(x).
\]
则当 \(t\to\infty\) 时，
\begin{equation}\label{eq:xi-poisson-cauchy-mixture-kl-t4-minimal-third}
D_{\mathrm{KL}}(p_t\mid q_t)
=\frac{\sigma^4}{8t^4}
+O\!\left(\frac{M_6}{t^6}\right).
\end{equation}
\end{theorem}

\begin{proof}
对固定 \(x\)，将 \(q_t(x-y)\) 在 \(y=0\) 处作四阶 Taylor 展开并对 \(\mu\) 积分。记中心矩 \(m_k:=\int y^k\,\dd\mu(y)\)（\(k\ge 2\)），则
\[
p_t(x)
=q_t(x)
+\frac{\sigma^2}{2}q_t''(x)
-\frac{m_3}{6}q_t^{(3)}(x)
+\frac{m_4}{24}q_t^{(4)}(x)
+R_t(x).
\]
由 Cauchy 核导数的有理闭式
\[
q_t^{(k)}(x)=\pi^{-1}t^{-(k+1)}R_k(x/t)\,(1+(x/t)^2)^{-(k+1)}
\]
可得统一余项控制 \(R_t(x)=O(M_6t^{-7})(1+(x/t)^2)^{-1}\)。因此存在有界有理函数 \(f_2,f_3,f_4\)（奇偶性分别为偶、奇、偶）使
\[
\frac{p_t(ty)}{q_t(ty)}-1
=\frac{\sigma^2}{t^2}f_2(y)
+\frac{m_3}{t^3}f_3(y)
+\frac{m_4}{t^4}f_4(y)
+O\!\left(\frac{M_6}{t^5}\right),
\]
其中
\[
f_2(y)=\frac{-1+3y^2}{(1+y^2)^2},
\qquad
\frac{q_t''(x)}{q_t(x)}=\frac{2(-t^2+3x^2)}{(t^2+x^2)^2}.
\]
记 \(\varepsilon_t:=p_t/q_t-1\)，则
\[
D_{\mathrm{KL}}(p_t\mid q_t)
=\int_{\RR}q_t(1+\varepsilon_t)\log(1+\varepsilon_t)\,\dd x.
\]
由 \(\int q_t\varepsilon_t\,\dd x=0\) 与
\(
(1+u)\log(1+u)=u+\frac12u^2+O(u^3)
\)
得
\[
D_{\mathrm{KL}}(p_t\mid q_t)
=\frac12\int_{\RR}q_t\varepsilon_t^2\,\dd x
+O\!\left(\int_{\RR}q_t|\varepsilon_t|^3\,\dd x\right).
\]
作代换 \(x=ty\) 并利用 \(q_t(ty)\,t\,\dd y=\pi^{-1}(1+y^2)^{-1}\dd y\)。由于权重为偶函数且 \(f_2\) 偶、\(f_3\) 奇，\(t^{-5}\) 的交叉项积分为零，故
\[
\frac12\int q_t\varepsilon_t^2\,\dd x
=\frac{\sigma^4}{2t^4}\cdot
\frac{1}{\pi}\int_{\RR}\frac{f_2(y)^2}{1+y^2}\,\dd y
+O\!\left(\frac{M_6}{t^6}\right).
\]
又
\[
\frac{f_2(y)^2}{1+y^2}
=\frac{1-6y^2+9y^4}{(1+y^2)^5},
\qquad
\frac{1}{\pi}\int_{\RR}\frac{f_2(y)^2}{1+y^2}\,\dd y=\frac14.
\]
此外 \(\int q_t|\varepsilon_t|^3\,\dd x=O(t^{-6})\)。合并即得
\eqref{eq:xi-poisson-cauchy-mixture-kl-t4-minimal-third}。证毕。
\end{proof}

\begin{theorem}[归一化熵能与耗散率极限]\label{cor:xi-poisson-cauchy-mixture-i1-t5-minimal-third}
\leanverified{paper\_xi\_poisson\_cauchy\_mixture\_i1\_t5\_minimal\_third}
在定理 \ref{thm:xi-poisson-cauchy-mixture-kl-t4-minimal-third} 的设定下，定义
\[
\mathcal E(t):=t^4D_{\mathrm{KL}}(p_t\mid q_t).
\]
则
\[
\lim_{t\to\infty}\mathcal E(t)=\frac{\sigma^4}{8},
\qquad
\lim_{t\to\infty}\left(-t^5\frac{\dd}{\dd t}D_{\mathrm{KL}}(p_t\mid q_t)\right)=\frac{\sigma^4}{2}.
\]
\end{theorem}

\begin{proof}
由 \eqref{eq:xi-poisson-cauchy-mixture-kl-t4-minimal-third},
\[
D_{\mathrm{KL}}(p_t\mid q_t)
=\frac{\sigma^4}{8t^4}+O(t^{-6}),
\]
乘以 \(t^4\) 得第一式。
上一条证明中的积分余项可逐项求导，故
\[
\frac{\dd}{\dd t}D_{\mathrm{KL}}(p_t\mid q_t)
=-\frac{\sigma^4}{2t^5}+O(t^{-7}),
\]
乘以 \(-t^5\) 并取极限即得第二式。证毕。
\end{proof}

\begin{corollary}[黄金分辨率链上的四阶衰减律]\label{cor:xi-poisson-cauchy-mixture-golden-exponent4-barrier}
\leanverified{paper\_xi\_poisson\_cauchy\_mixture\_golden\_exponent4\_barrier}
令
\[
\varphi:=\frac{1+\sqrt5}{2},\qquad
a_m:=m\log\varphi,\qquad
\varrho_m:=\tanh\!\left(\frac{a_m}{2}\right)\uparrow 1.
\]
定义 Cayley--Poisson 尺度
\[
t_m:=\frac{1+\varrho_m}{1-\varrho_m}.
\]
则 \(t_m=e^{a_m}=\varphi^m\)，并且
\[
D_{\mathrm{KL}}(p_{t_m}\mid q_{t_m})
=\frac{\sigma^4}{8\,\varphi^{4m}}
+O\!\left(\varphi^{-6m}\right)
\qquad(m\to\infty).
\]
\end{corollary}

\begin{proof}
恒等式 \(\frac{1+\tanh(x)}{1-\tanh(x)}=e^{2x}\) 给出 \(t_m=e^{a_m}=\varphi^m\)。
代入 \eqref{eq:xi-poisson-cauchy-mixture-kl-t4-minimal-third} 即得。
\end{proof}

\begin{theorem}[四阶尺度的刚性判别]\label{thm:xi-poisson-cauchy-mixture-t4-optimality-rigidity}
\leanverified{paper\_xi\_poisson\_cauchy\_mixture\_t4\_optimality\_rigidity}
在定理 \ref{thm:xi-poisson-cauchy-mixture-kl-t4-minimal-third} 的条件下，下列命题等价：
\begin{enumerate}
  \item \(D_{\mathrm{KL}}(p_t\mid q_t)=o(t^{-4})\)（\(t\to\infty\)）；
  \item \(\sigma^2=0\)；
  \item \(\mu=\delta_0\)（从而 \(p_t\equiv q_t\) 对所有 \(t>0\) 成立）。
\end{enumerate}
\end{theorem}

\begin{proof}
由 \eqref{eq:xi-poisson-cauchy-mixture-kl-t4-minimal-third},
\[
t^4D_{\mathrm{KL}}(p_t\mid q_t)=\frac{\sigma^4}{8}+O(t^{-2}),
\]
故 (1)\(\Rightarrow\)(2)。
又 \(\sigma^2=0\) 等价于 \(\int y^2\,\dd\mu(y)=0\)，由 \(y^2\ge0\) 得 \(\mu=\delta_0\)，即 (2)\(\Rightarrow\)(3)。
若 (3) 成立，则 \(p_t=q_t\) 且 KL 恒为零，故 (3)\(\Rightarrow\)(1)。证毕。
\end{proof}

\begin{corollary}[统一矩界与桥接接口下的条件推导]\label{thm:xi-poisson-kl-bridge-variance-necessary-condition}\label{cor:xi-poisson-kl-bridge-main-channel-dichotomy}
\leanverified{paper\_xi\_poisson\_kl\_bridge\_variance\_necessary\_condition}
设对每个 \(x_0\in\RR\) 给定中心化概率测度 \(\mu^{(x_0)}\)（\(\int y\,\dd\mu^{(x_0)}(y)=0\)），且存在统一矩界
\[
\sup_{x_0}\int y^2\,\dd\mu^{(x_0)}(y)\le \Sigma^2,
\qquad
\sup_{x_0}\int |y|^6\,\dd\mu^{(x_0)}(y)\le M_6<\infty.
\]
令
\[
p_t^{(x_0)}:=P_t*\mu^{(x_0)},\qquad q_t:=P_t.
\]
则
\[
\sup_{x_0\in\RR}D_{\mathrm{KL}}(p_t^{(x_0)}\mid q_t)
\le
\frac{\Sigma^4}{8t^4}
+O\!\left(\frac{M_6}{t^6}\right),
\]
并沿 \(t_m=\varphi^m\) 有
\[
\sup_{x_0}D_{\mathrm{KL}}(p_{t_m}^{(x_0)}\mid q_{t_m})
\le
\frac{\Sigma^4}{8\varphi^{4m}}
+O(\varphi^{-6m}).
\]
进一步，若存在常数 \(C>0\) 使
\[
\sup_{x_0\in\RR}D_{\zeta,x_0}(\varrho)
\le
C\sup_{x_0\in\RR}D_{\mathrm{KL}}\!\bigl(p_{t(\varrho)}^{(x_0)}\mid q_{t(\varrho)}\bigr)
+o(1)
\qquad(\varrho\uparrow1),
\]
其中 \(t(\varrho):=\frac{1+\varrho}{1-\varrho}\)，则沿 \(\varrho=\varrho_m\) 有
\[
\sup_{x_0}D_{\zeta,x_0}(\varrho_m)\to0.
\]
于是由定理 \ref{thm:xi-rh-iff-uniform-comoving-defect} 可推出 \(\mathrm{RH}\)。
\end{corollary}

\begin{proof}
对每个 \(x_0\) 应用定理 \ref{thm:xi-poisson-cauchy-mixture-kl-t4-minimal-third}，并用统一矩界控制主项与余项常数，即得第一组估计。
取 \(t=t_m=\varphi^m\) 即得第二组估计。
若桥接不等式成立，则右端在 \(\varrho_m\uparrow1\) 时趋于零，从而 \(\sup_{x_0}D_{\zeta,x_0}(\varrho_m)\to0\)。
再由定理 \ref{thm:xi-rh-iff-uniform-comoving-defect} 的等价判据即得结论。
\end{proof}

\endinput
